% !TeX root = ../main.tex
\documentclass[../main]{subfiles}
\begin{document}
\section*{Programa de estudio de la asignatura dibujo técnico asistido}
\begin{multicols}{2}
  \begin{enumerate}
    \item BLOQUE 1
      \begin{itemize}
        \item Letras, formato de láminas.
        \item Líneas aplicaciones.
        \item Interrupción de cuerpos.
        \item Líneas de rotura, aplicaciones.
        \item Rayados y colores convencionales.
        \item Modo de plegar planos.
        \item Rotulación de planos.
      \end{itemize}
    \item BLOQUE 2
      \begin{itemize}
        \item Métodos de representación:sistema IRAM.\@
        \item Sistemas de representación: perspectiva caballera,
          perspectiva isométrica.
        \item Rebatimiento y desarrollos.
      \end{itemize}
    \item BLOQUE 3
      \begin{itemize}
        \item Acotaciones.
        \item Acotaciones de los dibujos.
        \item Clasificación de cotas.
        \item Ejecución de líneas de cota y referencias.
        \item Sistemas de acotamiento.
        \item Acotaciones condicionadas a líneas base o de referencia.
      \end{itemize}
    \item BLOQUE 4
      \begin{itemize}
        \item Generalidades del software Autocad.
        \item Estructura de pantalla.
        \item Símbolos UCS.\@
        \item Modos de seleccionar objetos.
        \item Modos de captura de objetos.
        \item Comandos: línea, rectángulos, polígonos, copiar, pegar,
          recortar, rotar, espejar, trazar paralelas, empalmar líneas,
          borrar, extender una línea, escribir un texto, distancia,
          área, arcos, circunferencias.
        \item Editar un texto.
        \item Cargar distintos tipos de línea.
        \item Generar bloques.
        \item Insertar bloques.
        \item Nociones sobre manejo de capas, formas de acotar,
          impresión de trabajos.
      \end{itemize}
  \end{enumerate}
\end{multicols}
\end{document}
